\documentclass[11pt,a4paper]{article}
\usepackage{series}
\usepackage{booktabs}
\usepackage{array}
\usepackage{xcolor}

\newcommand{\End}{\operatorname{End}}
\newcommand{\Aut}{\operatorname{Aut}}
\newcommand{\Sym}{\operatorname{Sym}}
\newcommand{\Gr}{\operatorname{Gr}}
\newcommand{\QFT}{\mathrm{QFT}}
\newcommand{\CAR}{\mathrm{CAR}}
\newcommand{\CCR}{\mathrm{CCR}}
\newcommand{\BV}{\mathrm{BV}}
\newcommand{\BRST}{\mathrm{BRST}}
\newcommand{\MTT}{\mathrm{MTT}}

\title{\textbf{Universal Perturbative Graph Structure:}\\
\large A Typed Modal Diagrammatics and Its Quantum Boundary}
\author{Peter Nero}
\date{Version 3, September 2026}

\begin{document}
\maketitle

\begin{abstract}
Feynman graphs are used in quantum field theory, classical statistical
field theory, stochastic dynamics, and effective descriptions of gravity.
Their recurrence has a precise but limited explanation.  Once one specifies
a nondegenerate quadratic kernel, interaction tensors, a grading, and a
contraction rule, Wick expansion organizes perturbative coefficients by
graphs.  The graph grammar is therefore not uniquely quantum.  It also does
not, by itself, select a physical propagator, Bose or Fermi statistics,
causal locality, gauge reduction, counterterms, positivity, or unitarity.

This paper reformulates modal diagrammatics around that distinction.  We
prove a finite graded graph-expansion theorem and an exact projection
theorem: if an even projector intertwines the quadratic operator,
interaction tensors, grading, and declared identities, then every
all-retained graph coefficient agrees with the coefficient computed in the
projected theory.  We also prove a nonselection theorem showing that the
same graph combinatorics supports continuously many inequivalent weights
and different statistics.  Thus basin adjacency or coherence language
alone cannot derive Feynman rules.

For Modal Triplet Theory, the result supplies a rigorous transfer target.
The selected q79 twisted-Dirac construction already composes with standard
CAR/AQFT machinery to give a free even local net.  The present theorem
explains how a selected upper action could transfer its perturbative graph
data to a lower effective sector.  It does not manufacture that action.
The geometry-selected upper action, fixed-coupling interacting
gauge--BRST \(C^*\)-completion, renormalization-group matching, and
observable uncertainty packet remain open.  The formulation
retains the explanatory value of modal diagrammatics while making its
mathematical and physical boundary explicit. A finite residual-and-metric
jet supplies the complete repair-action graph datum, with exact transport
under a compatible change of variables; its identification with physical
Feynman data remains a separate action and state problem.
\end{abstract}

\section*{Version 3 revision note}
\begin{description}
\item[Supersedes.] Version~2; the earlier revision note and released identity
are retained.
\item[Reason.] The residual-to-graph construction now supplies an upstream
input to the transfer theorem. The graph formula also needed an explicit
interaction sign and correctly typed field/dual projection.
\item[Resolution.] Explain the complete residual/metric jet, a polynomial
example, and non-isometric coordinate transport. Distinguish cost pullback
from reducing isometry and include the Euclidean vertex sign.
\item[Retained result.] Finite Wick organization, all-retained graph
transfer, no-mixed-vertex factorization, and graph nonselection remain exact.
\item[Open boundary.] A positive repair cost does not select the physical
signed action, causal inverse, statistics, state, or interacting completion.
\end{description}

\section*{Version 2 revision note}

\paragraph{Supersedes.} Version 1.0 (DOI
\href{https://doi.org/10.5281/zenodo.18330804}{10.5281/zenodo.18330804}).

\paragraph{Reason.} Version 1 correctly emphasized that diagrammatic expansions are not
exclusive to quantum theories, but it attributed too much to coherence and
Gaussian combinatorics.  In several places it treated a modal Hessian as a
selected physical propagator, described projected diagrams as necessarily
the standard Feynman rules, and discussed quantum behavior and gravity
beyond what the supplied structure proved.

\paragraph{Resolution.} Version 2 makes four substantive changes.
\begin{enumerate}
\item It replaces the broad origin claim by exact finite graded
graph-expansion and projection theorems.
\item It distinguishes the quadratic inverse, Euclidean covariance,
retarded/advanced Green operators, causal propagator, state two-point
function, and Feynman kernel.
\item It proves that graph incidence data do not select numerical weights,
statistics, gauge identities, counterterms, or a quantum theory.
\item It reconciles the paper with the current MTT ledger: the selected
free q79 even-CAR net is closed at its declared tier, while the upper action
and nonperturbative interacting completion remain open.
\end{enumerate}
These are corrections to the central claim, not merely added caveats.

\paragraph{Retained result.} The original observation that classical and
quantum models can share the same perturbative graph grammar is retained
and is now proved at the exact finite graded level.

\paragraph{Remaining boundary.} A selected upper action, its complete
projection intertwiners, and a fixed-coupling interacting gauge--BRST
\(C^*\)-completion are not established here.

\section{Why the same diagrams recur}

A perturbative graph is a compact record of two operations:
\begin{enumerate}
\item insert multilinear interaction tensors;
\item contract their arguments pairwise with a chosen two-point kernel.
\end{enumerate}
The record is useful because the number of pairings grows rapidly.
Wick's theorem is the algebraic mechanism that turns those pairings into
graph sums \cite{Wick1950}.  Nothing in that combinatorial statement alone
decides whether the underlying system is a Euclidean random field, a
Lorentzian quantum field, a response theory, or a finite algebraic model.

This observation motivates, but does not yet prove, the MTT interpretation.
A fixed point can provide a background about which to expand.  A projector
can identify retained modes.  Neither object determines the action, its
quadratic inverse, or its higher derivatives.  Those data must be supplied
or derived separately.  The useful question is therefore not
\emph{whether coherence automatically creates Feynman rules}, but:
\begin{quote}
Under which explicit hypotheses does perturbative graph data pass from an
upper description to a projected effective description?
\end{quote}

The answer below is exact on a declared finite or regularized domain.  It
is intentionally not a claim about convergence of an infinite
perturbation series or existence of an interacting continuum QFT.

\section{The typed perturbative input}

\subsection{Finite graded field space}

Let
\[
 E=E_{\bar 0}\oplus E_{\bar 1}
\]
be a finite-dimensional real or complex super vector space.  Even
coordinates represent bosonic directions and odd coordinates represent
fermionic directions.  Fix:
\begin{enumerate}
\item an even nondegenerate quadratic operator \(K:E\to E^*\);
\item its inverse contraction kernel \(G=K^{-1}:E^*\to E\);
\item graded-symmetric interaction tensors
\[
 V_n\in (E^*)^{\otimes n},\qquad 3\leq n\leq N;
\]
\item a formal coupling parameter \(\lambda\);
\item an ordering and orientation convention for odd half-edges.
\end{enumerate}
The formal action is
\begin{equation}
 S(\phi)=\frac12 K(\phi,\phi)
 +\sum_{n=3}^{N}\frac{\lambda^{n-2}}{n!}V_n(\phi,\ldots,\phi).
 \label{eq:action}
\end{equation}
The powers of \(\lambda\) are conventional and may be replaced by
independent formal couplings without changing the argument.

The word ``typed'' matters.  The pair \((K,V_n)\) does not yet define a
physical QFT.  It defines algebraic graph weights.  A physical
interpretation requires additional structures described in
\cref{sec:quantum-boundary}.

\subsection{Graph weights}

Let \(\Gamma\) be a finite graph whose vertices have valence between \(3\)
and \(N\).  Assign \(V_{\deg(v)}\) to every vertex and \(G\) to every
internal edge.  Contract matching tensor indices.  For odd labels, include
the Koszul sign \(\epsilon(\Gamma)\) determined by the declared ordering.
The resulting scalar or external tensor is
\begin{equation}
 W_\Gamma(K,V,G)
 =
 \frac{(-1)^{|V(\Gamma)|}\epsilon(\Gamma)}{|\Aut\Gamma|}
 \operatorname{Contr}_\Gamma
 \left(
 \bigotimes_{v\in V(\Gamma)}V_{\deg(v)}
 \otimes
 \bigotimes_{e\in E_{\mathrm{int}}(\Gamma)}G
 \right).
 \label{eq:graph-weight}
\end{equation}
This formula exposes all of the information that a bare graph drawing
hides: numerical kernels, vertex tensors, symmetry factors, and grading.
The displayed sign is for the Euclidean convention
$\exp(-S_{\rm int})$ with $\hbar=1$; coupling powers are understood.
Other exponential conventions replace that vertex factor explicitly.
External insertions are labeled and graph automorphisms fix those labels.
This formula first organizes unnormalized moments; division by the vacuum
series gives normalized correlators and removes disconnected vacuum factors.

\begin{definition}[Perturbative graph datum]
A perturbative graph datum is the tuple
\[
 \mathfrak D=(E,\Gr E,K,G,\{V_n\},\epsilon,\mathcal R),
\]
where \(\mathcal R\) records any regularization or renormalization map
needed to define singular contractions.  In a finite nonsingular model,
\(\mathcal R=\Id\).
\end{definition}

\section{The graph expansion theorem}

\begin{theorem}[Finite graded Wick graph expansion]
\label{thm:wick-graph}
Let \(E\), \(K\), and \(V_n\) be as above.  Treat the Gaussian expectation
with covariance \(G=K^{-1}\) and the interaction exponential as a formal
power series.  Every coefficient of every polynomial correlation
function is a finite sum of the weights \eqref{eq:graph-weight}.  The
vertices are the tensors \(V_n\), the internal edges are contractions by
\(G\), the automorphism denominator gives the symmetry factor, and the
odd permutation convention gives the fermionic sign.
\end{theorem}

\begin{proof}
Expand the interaction exponential order by order in \(\lambda\).  At a
fixed order only finitely many products of the \(V_n\) occur.  Apply the
graded Wick rule to each polynomial Gaussian expectation.  Every complete
pairing identifies two tensor arguments and contracts them with \(G\).
Grouping pairings that differ only by relabeling yields an isomorphism
class of graphs and the orbit--stabilizer factor
\(|\Aut\Gamma|^{-1}\).  Moving odd variables into the chosen pairing order
produces the Koszul sign, while the interaction exponential contributes
one minus sign per vertex. Conversely, every decorated graph determines
such a pairing.  Thus the two finite sums agree coefficient by
coefficient.
\end{proof}

\begin{remark}
The theorem is exact as a statement about formal coefficients.  It does
not assert convergence, reflection positivity, causal factorization,
unitarity, or existence of a continuum measure.
\end{remark}

\subsection{What is universal}

The universal content is now precise:
\[
\begin{gathered}
\text{quadratic contraction}
+\text{multilinear interactions}
+\text{grading}\\
\Longrightarrow\quad
\text{graph-indexed formal coefficients}.
\end{gathered}
\]
Classical Gaussian measures and quasifree quantum states both instantiate
this schema.  Their physical meanings remain different because their
algebras, kernels, state conditions, and localization rules differ.

\section{Exact transfer through a finite projection}

\subsection{Compatible retained sector}

Let \(P:E\to E\) be an even idempotent and write \(E_P=\im P\).
Let $P^*:E^*\to E^*$ be its dual. Because $K:E\to E^*$ is a
quadratic form, the correctly typed block-diagonality condition is
\begin{equation}
 KP=P^*K.
 \label{eq:PK}
\end{equation}
Then \(K\) is block diagonal on \(E_P\oplus\ker P\), and
\[
 K_P=K|_{E_P},\qquad G_P=K_P^{-1}=PGP^*|_{E_P^*}.
\]
Here $E_P^*$ is embedded by the splitting determined by $P$. If a field
metric identifies $E$ with $E^*$ and $P$ is orthogonal for that metric,
this condition becomes commutation of $P$ with the Riesz operator of $K$.
Define the retained vertex
\[
 V_n^P=V_n|_{E_P^{\otimes n}}.
\]

\begin{theorem}[All-retained graph transfer]
\label{thm:projection}
Under \eqref{eq:PK}, let \(\Gamma\) be a decorated graph for which every
half-edge label lies in \(E_P\).  Then
\[
 W_\Gamma(K_P,V^P,G_P)
 =
 W_\Gamma(K,V,G)\big|_{E_P}.
\]
The equality includes symmetry factors and graded signs.
\end{theorem}

\begin{proof}
At every retained vertex, \(V_n^P\) is exactly the restriction of \(V_n\).
On every retained internal edge, commutation and block diagonality give
\(G_P=PGP^*|_{E_P^*}\). Since \(P\) is even, it preserves parity and hence
the Koszul sign.  The underlying graph and its automorphism group are
unchanged.  Substitution into \eqref{eq:graph-weight} proves the equality.
\end{proof}

\begin{corollary}[Closed retained subtheory]
\label{cor:closed-sector}
If, in addition, every mixed interaction component vanishes,
\[
 V_n(x_1,\ldots,x_n)=0
 \quad\text{whenever retained and discarded arguments both occur},
 \label{eq:no-mixed}
\]
then the full formal perturbative expansion with retained external legs
factorizes from the discarded sector.  After normalization by vacuum
terms, it equals the expansion computed solely from
\((E_P,K_P,\{V_n^P\})\).
\end{corollary}

\begin{proof}
Condition \eqref{eq:no-mixed} prevents a connected graph with retained
external legs from containing a discarded internal label.  Apply
\cref{thm:projection} to every connected contribution.  Purely discarded
vacuum graphs factor and cancel in normalized correlators.
\end{proof}

\subsection{Why effective actions are needed}

Without \eqref{eq:no-mixed}, the theorem does \emph{not} say that simply
deleting discarded modes gives the correct lower theory.  Mixed vertices
allow discarded internal lines to contribute to retained observables.
Integrating them out produces an effective action
\[
 e^{-S_{\mathrm{eff}}(\phi_P)}
 =
 \int_{\ker P}e^{-S(\phi_P+\chi)}\,\dd\chi
\]
in a Euclidean finite model, or its appropriate formal/algebraic analogue.
That effective action usually contains new nonlocal or higher-order
vertices.  Projection and marginalization are therefore different
operations.

\subsection{Transfer of identities}

\begin{proposition}[Intertwined identity transfer]
\label{prop:ward}
Let \(\mathcal A\) and \(\mathcal A_P\) be algebras of formal functionals,
let \(R:\mathcal A\to\mathcal A_P\) be the restriction induced by \(P\),
and let \(\mathcal W,\mathcal W_P\) be linear identity operators.  If
\[
 R\mathcal W=\mathcal W_P R,
 \]
then \(\mathcal W F=0\) implies \(\mathcal W_P(RF)=0\).
\end{proposition}

\begin{proof}
Apply \(R\) to \(\mathcal W F=0\) and use the intertwining equality.
\end{proof}

This proposition is elementary, but it identifies the real obligation in
gauge theory.  Ward, BRST, or BV identities descend only after their
operators and brackets are shown to intertwine.  A mode count or a graph
match does not establish this.

\section{The nonselection theorem}

\begin{theorem}[Graph combinatorics does not select Feynman rules]
\label{thm:nonselection}
The incidence structure of perturbative graphs, even together with a list
of allowed valences, does not determine:
\begin{enumerate}
\item propagator values;
\item interaction strengths;
\item bosonic or fermionic statistics;
\item causal or boundary prescriptions;
\item gauge-fixing, ghost, or BRST data;
\item renormalization extensions or counterterms.
\end{enumerate}
In particular, continuously many inequivalent perturbative graph data have
the same underlying graph set.
\end{theorem}

\begin{proof}
Take the one-dimensional bosonic action
\[
 S_{k,g}(x)=\frac{k}{2}x^2+\frac{g}{4!}x^4,
 \qquad k>0.
\]
For every \(k\) and \(g\neq0\), the allowed connected graphs are the same
four-valent graphs.  A graph with \(I\) internal edges and \(V\) vertices
has weight proportional to \(k^{-I}g^V\).  Varying \(k\) or \(g\)
therefore changes all nontrivial weights without changing incidence data.
Replacing the field by an odd multiplet changes exchange signs while the
ungraded incidence graph can remain identical.  Retarded, advanced,
Euclidean, and Feynman prescriptions may solve equations associated with
the same differential expression but obey different support or boundary
conditions.  Gauge fixing adds ghost and identity data, and
renormalization chooses extensions of singular distributions.  None is
encoded by the bare graph set.
\end{proof}

\begin{corollary}[Basin adjacency is insufficient]
A basin graph, neighborhood relation, or coherence-capacity network cannot
by itself derive physical Feynman rules.  It must be supplemented by a
selected quadratic operator and inverse prescription, interaction tensors,
grading, gauge structure where applicable, and a renormalization rule.
\end{corollary}

\section{Which kernel is meant?}
\label{sec:quantum-boundary}

Version 1 used ``propagator'' too broadly.  The following objects can be
related, but they are not interchangeable.

\begin{center}
\small
\begin{tabular}{>{\raggedright\arraybackslash}p{0.23\textwidth}
                >{\raggedright\arraybackslash}p{0.31\textwidth}
                >{\raggedright\arraybackslash}p{0.35\textwidth}}
\toprule
Object & Defining input & What it does not automatically provide\\
\midrule
Algebraic inverse \(K^{-1}\) &
Nondegenerate finite quadratic form &
Causality, positivity, or a quantum state\\
Euclidean covariance &
Positive elliptic operator and measure &
Lorentzian time ordering without reconstruction hypotheses\\
Retarded/advanced Green operator &
Green-hyperbolic operator and support condition &
A state two-point function\\
Causal propagator &
Difference of retarded and advanced operators &
Positive-frequency or Hadamard choice\\
Two-point function &
Observable algebra and state &
Time ordering or unique vacuum\\
Feynman kernel &
Time ordering/state or boundary prescription &
Nonperturbative completion\\
\bottomrule
\end{tabular}
\end{center}

On curved spacetime, local covariant Wick products and time-ordered
products require microlocal and renormalization conditions
\cite{BrunettiFredenhagen2000,HollandsWald2001,HollandsWald2002}.
Epstein--Glaser renormalization shows that locality and causal
factorization constrain the extension of distributions
\cite{EpsteinGlaser1973}.  These are physical and analytic inputs, not
consequences of graph drawing.

\section{Gauge theories and renormalization}

\subsection{Gauge-fixed perturbation theory}

For a gauge theory the naive Hessian is degenerate along gauge orbits.
One must introduce gauge-fixing and ghost sectors, or work in a
homological BV/BRST formulation.  The interaction data then include
graded fields, antifields, a bracket, and a differential.  The quantum
master equation or corresponding Ward identities control the consistency
of the perturbative expansion \cite{Rejzner2016}.

Modal projection is compatible with such a construction only if the
projection preserves the relevant complex and identities.  Schematically,
\[
 PQ=QP,\qquad
 P\{\cdot,\cdot\}=\{P\cdot,P\cdot\}_P,
\qquad
 R\mathcal W=\mathcal W_P R.
\]
These equations are proof obligations.  Calling discarded directions
``inadmissible'' does not make them true.

\subsection{Renormalized graph weights}

In a continuum theory, products of propagators can be singular on
coincident-point diagonals.  A graph before extension may define a
distribution only away from those diagonals.  Renormalization extends it
subject to locality, covariance, scaling, and symmetry requirements.
Different allowed extensions correspond to finite renormalization
freedoms.  Thus the renormalization map \(\mathcal R\) belongs in the
perturbative graph datum.

The graph theorem remains useful: it organizes which distribution must be
extended.  It does not perform or uniquely select the extension.

\section{Repair jets supply a finite graph datum}
\label{sec:repair-jets}

The frozen H4-T8 construction supplies an upstream arrow to this paper's
graph theorem \cite{FrozenRepairJet}. Fix finite field and residual spaces,
a background $a_*$, a smooth or formal residual $\Phi(a)$, and a positive
residual metric field $W(a)$. Set
\[
 S_{\rm rep}(a)=\tfrac12\langle\Phi(a),W(a)\Phi(a)\rangle.
\]
Write $R_m=D^m\Phi(a_*)$ and $W_m=D^mW(a_*)$. Repeated Leibniz
differentiation assigns each derivative to the left residual, metric, or
right residual, so the entire vertex tensor is
\[
 V_n(v_1,\ldots,v_n)=\tfrac12
 \sum_{[n]=A\sqcup B\sqcup C}
 \langle R_{|A|}(v_A),W_{|B|}(v_B)R_{|C|}(v_C)\rangle.
\]
The blocks are ordered by their factor roles and indices retain their
original order; graded models insert the corresponding Koszul signs.
Thus vertices are outputs of the joint residual-and-metric jet, not extra
source coefficients. A residual without its metric derivatives is insufficient.
At $R_0=0$, the Hessian form is $K=R_1^\dagger W_0R_1$.
At nonzero defect even constant $W$ adds
$\langle R_0,W_0R_2(v,w)\rangle$, so the Gram formula is no longer complete.

\subsection{A worked polynomial source}
\label{sec:repair-example}

For $\Phi(x,y)=(x,y+y^2)$ and $W=I$ at zero,
\[
 S_{\rm rep}=\tfrac12x^2+\tfrac12y^2+y^3+\tfrac12y^4,
 \qquad K=I,\quad V_{3,yyy}=6,\quad V_{4,yyyy}=12.
\]
All higher vertices vanish. In general a residual of degree $r$ with
constant metric has no vertices above degree $2r$. With the same scalar
residual and $W(y)=1+y+y^2>0$, the cubic and quartic derivatives instead
become $9$ and $48$, and the action has degree six. This is a metric-jet
effect, not new independent interaction data.

In the unit-covariance Gaussian, $\langle y^4\rangle=3$ and
$\langle y^6\rangle=15$. The unnormalized vacuum coefficients from one
quartic vertex and two cubic vertices are consequently
$-(12/4!)3=-3/2$ and $(6/3!)^2 15/2=15/2$.
These illustrate the exponential sign as well as the combinatorics.

\subsection{Coordinate transport is not a physical isometry}
\label{sec:repair-transport}

If linear bijections $T,U$ satisfy
$\Phi'(Ta)=U\Phi(a)$ and
$\langle Ur,W'(Ta)Us\rangle=\langle r,W(a)s\rangle$,
then $S'_{\rm rep}(Ta)=S_{\rm rep}(a)$ and every vertex pulls back.
For $T=\operatorname{diag}(2,3)$ in the example, the target Hessian
form is $\operatorname{diag}(1/4,1/9)$, its contravariant Gaussian
covariance is $\operatorname{diag}(4,9)$, and its cubic/quartic entries
are $2/9,4/27$. Edge and vertex factors cancel exactly in every contraction.
Transporting the field metric also makes the Hessian Riesz operators and
their functional calculi intertwine. With independently fixed physical
metrics, a cost pullback alone does not prove that stronger statement.

Nor is a lower-dimensional cost restriction a reducing embedding. For
$K=\left(\begin{smallmatrix}1&1\\1&2\end{smallmatrix}\right)$ and
$Tx=(x,0)$, the restricted Hessian is $T^*KT=1$, but
$KT(1)=(1,1)\ne T(1)$. The retained line leaks into its complement.
The all-retained theorem therefore still needs~\eqref{eq:PK}, and the
full retained theory still needs the no-mixed-vertex condition or an
effective-action calculation.

\subsection{What the source reduction does not quantize}

On a specified kernel complement, the repair Hessian yields a finite
Euclidean reduced inverse, zero-mode projector, and heat semigroup.
These are sufficient for formal repair graphs. The joint q79 residual can
organize integrability, HYM, balance, anomaly, and normalization lanes,
but its actual endpoint and metric jets are not evaluated by H4-T8.
In particular the signed-action comparison in the action paper
\cite{ClosureAction5} is still necessary: squaring an action gradient
changes both its Hessian and its vertices. Neither a causal/Feynman
prescription nor a physical gauge/BRST state is obtained by differentiating
a positive norm. The selected free-CAR and operational-QM results remain
at their established tiers; this construction does not reopen them or
complete the interacting theory.

\section{What MTT contributes now}

\subsection{Fixed points and coherent sectors}

MTT supplies a language of coherent fixed points, admissible sectors, and
finite projections.  At the level needed here, a fixed point identifies a
background \(u_*\), and a selected action \(S\) would yield
\[
 K=\dd^2S[u_*],
 \qquad
 V_n=\dd^nS[u_*].
\]
If \(S\), \(u_*\), and the inverse prescription are selected from one
source, the tensors in the graph datum are no longer arbitrary.  If only
the fixed point or carrier is known, they remain underdetermined.

This is the significance of the current upper-action blocker.  The missing
theorem is not another restatement of Wick expansion.  It is a same-source
construction of an upper differential/action whose Hessian, higher
products, symmetries, and projections reproduce the accepted lower
structures.

\subsection{The selected free q79 CAR result}

The companion QFT paper owns the current selected free-field theorem
\cite{MTTQFT2026}.  On the declared globally hyperbolic framed q79
representative, the selected twisted massless Dirac operator composes with
standard Green-hyperbolic and CAR/AQFT machinery to produce an even local
observable net with locality, covariance, the time-slice property, and a
nonempty positive Hadamard state space.

The present paper does not re-prove that theorem.  It uses it to mark the
quantum boundary accurately:
\[
\text{selected Dirac source}
\longrightarrow
\text{standard CAR/AQFT construction}
\longrightarrow
\text{free local quantum net}.
\]
The CAR relations and Hadamard machinery are not inferred from modal
graphs.

\subsection{The interacting boundary}

Current MTT work also records a classical BV master action, a
gauge-fixed Green-hyperbolic equicausal algebra, and a formal all-orders
anomaly-free quantum master equation at the declared perturbative tier.
Those results support a conditional perturbative graph construction.  A
formal jet at zero coupling does not select a unique fixed-nonzero-coupling
\(C^*\)-completion.  A geometry-selected nonperturbative gauge--BRST
bridge or regulator/continuum limit, selected physical state, RG matching,
uncertainty control, and observable comparison remain open.

\subsection{Finite transferred operations and the physical endpoint}
\label{sec:finite-transfer-consumer}

The cohesive source is a twisted perfect-object benchmark with integrable
superconnection \(\overline E\). Conjugation cancels its scalar overlap
twist on endomorphisms, giving the ordinary differential
\(d=[\overline E,-]\). A chosen Hermitian metric gives total-space Hodge
theory; it does not give a uniformly bounded family of fiberwise Green
operators through a cohomology-rank jump \cite{FrozenCohesive,Cohesive14}.
Its Maurer--Cartan residual \(da+a^2\), together with the linear gauge row,
has positive normal \(\Delta_1\) at the exact background. The resulting
heat semigroup is the tangent repair evolution, not the full nonlinear flow
or a signed physical action \cite{FrozenCohesiveMC}. A derived equivalence
preserving formal deformations does not automatically preserve that metric,
adjoint, or spectrum. This is why the repair-jet graph construction still
requires its own declared metric and action interpretation.

Cohesive's finite response hierarchy supplies a related but different
structure from the scalar repair jets above. Its 144-dimensional Weyl DGA
retracts onto a 48-dimensional response complex. A DGA morphism preserving
the inclusion, projection, and homotopy transports every transferred
\(m_n\); the proof moves that morphism through each decorated transfer tree.
The finite translation and Fourier actions preserve those source identities,
so their order-36 covariance acts at all arities
\cite{Cohesive14,FrozenAllArity}. This is not an inference from checking a few
vertices. Nor are the \(m_n\) already a signed physical action or measured
couplings. An approximate or nonreducing comparison needs separate defect
estimates and does not inherit exact Hodge or all-arity identities.

The endpoint-factorization theorem explains what physical promotion would
require \cite{FrozenEndpoint}. Geometry and signed action (GAS), spectral
synthesis and its domain/contraction certificates (SYN), and a
BV-compatible four-dimensional compactification (BV4) must share one source.
Their symmetry and Galerkin/Feshbach rows are then derived from that source,
not additional freely selected kernels. For example, fixing a retained
quadratic block does not fix the complement contribution to its Schur
complement. The frozen contract accepts zero of the three complete physical
packets and zero of its seven endpoint rows. These are structured source
obligations, not a three-parameter count. Their openness leaves both the
finite graph theorem and the selected free CAR net intact.

\section{Three illustrative cases}

\subsection{Euclidean scalar model}

For a positive matrix \(K\), the finite Gaussian reference measure is
literal. A quartic interaction defines a normalized probability measure only
when its full Boltzmann weight is integrable; a nonnegative quartic form is
one sufficient condition. For an arbitrary quartic tensor the Wick
coefficients still make sense formally, but positivity of \(K\) alone does
not make the interacting integral finite. This example shows that the graph
grammar is not uniquely quantum without equating a formal series with a
convergent integral.

\subsection{Free fermions}

For an odd field, the Gaussian integral is algebraic in Grassmann
variables, and contractions carry signs.  The ungraded graph shape alone
cannot recover those signs.  One needs the parity and ordering data.  In a
physical CAR theory one additionally needs the CAR algebra, localization,
and a state.

\subsection{Gauge fields}

For gauge fields, the uncorrected Hessian has a kernel along gauge
directions.  A propagator appears only after gauge-fixing or homological
reduction, and ghost vertices are part of the graph datum.  BRST or BV
identities constrain admissible counterterms.  This case makes especially
clear why ``inverse Hessian plus coherence'' is insufficient.

\section{Classical, quantum, and gravitational uses}

\subsection{Classical statistical fields}

In classical statistical field theory, a positive covariance and
interaction potential generate graph expansions of moments.  This is a
fully adequate explanation for the recurrence of the same combinatorial
grammar.  It does not imply that classical and quantum theories differ
only in interpretation: noncommutative observables, interference phases,
local commutation relations, and state conditions are additional
mathematical structures.

\subsection{Quantum field theory}

In QFT, Wick expansion acts on a quasifree state or time-ordered product.
The propagator prescription, statistics, local algebra, and
renormalization conditions give the graphs their physical meaning.  The
current MTT free-CAR theorem provides such a free quantum source on its
declared branch.  Interacting QFT requires the remaining bridge.

\subsection{Perturbative gravity}

Expanding a specified gravitational action around a background also
produces quadratic and higher tensors, so graph organization applies.
This neither proves UV completion nor shows that spacetime quantization is
unnecessary.  It establishes only the same conditional perturbative
grammar.  Any stronger MTT gravity claim must derive the background,
action, gauge complex, state, and ultraviolet control from the selected
source.

\section{Status and theorem ownership}

\begin{center}
\small
\begin{tabular}{>{\raggedright\arraybackslash}p{0.34\textwidth}
                >{\raggedright\arraybackslash}p{0.18\textwidth}
                >{\raggedright\arraybackslash}p{0.36\textwidth}}
\toprule
Statement & Status & Owner or boundary\\
\midrule
Finite graded Wick graph expansion &
Exact here &
\Cref{thm:wick-graph}\\
All-retained finite projection equality &
Exact here &
\Cref{thm:projection}\\
Closed-sector factorization under no mixed vertices &
Exact here &
\Cref{cor:closed-sector}\\
Graph-combinatorics nonselection &
Exact here &
\Cref{thm:nonselection}\\
Selected q79 free even-CAR net &
Closed at declared tier &
Companion MTT-to-QFT paper; imported CAR/AQFT theorems\\
Embedded renormalized-SM equivalence &
Closed at profile tier &
Current A04 result; not a derivation of quantization\\
Upper action and automorphism transfer &
Open &
Current blocker B.ACTION.01\\
Nonperturbative interacting gauge--BRST \(C^*\) bridge &
Open &
Current blocker B.QFT.02\\
No-knob Standard Model values &
Open &
Current strict-upgrade ledger\\
\bottomrule
\end{tabular}
\end{center}

This ownership table prevents a common form of accidental duplication.
The paper proves the finite graph and transfer statements.  It cites the
free-CAR result.  It does not promote formal or profile evidence into an
interacting-QFT theorem.

\section{Validation program}

A candidate upper-to-lower diagrammatic derivation should be accepted only
after the following checks.
\begin{enumerate}
\item \textbf{Source check:} identify the action, background, grading, and
all continuous inputs by path and hash.
\item \textbf{Quadratic check:} prove invertibility on the declared
gauge-fixed or reduced domain and specify the inverse prescription.
\item \textbf{Vertex check:} compute every retained \(V_n\) from the same
action, including overlap integrals and normalization conventions.
\item \textbf{Projection check:} verify \(KP=P^*K\), the vertex
intertwiners, and whether mixed vertices vanish or are integrated out.
\item \textbf{Statistics check:} preserve the \(\mathbb Z_2\) grading and
Koszul signs.
\item \textbf{Gauge check:} verify BRST/BV chain maps and Ward or QME
identities.
\item \textbf{Renormalization check:} specify distribution extensions,
scheme transport, and residual finite freedoms.
\item \textbf{Physical check:} establish local algebra, positivity/state
conditions, uncertainty control, and observable comparison at the claimed
tier.
\end{enumerate}
Passing only the graph-shape check is not enough.

\section{Discussion}

\subsection{What survives from the original idea}

The central intuition survives in a stronger form.  Feynman-style diagrams
are not mysterious pictures added to quantum mechanics.  They are the
natural combinatorial language of perturbative contractions.  MTT can use
that universality because its fixed points and projections offer a
candidate upstream organization of backgrounds and retained sectors.

The correction is equally important.  ``Universal graph grammar'' is not
``universal physical dynamics.''  A graph does not know which kernel is
causal, which state is positive, which field is fermionic, which gauge
identity must hold, or which counterterm is allowed.  Those facts reside
in the typed data attached to the graph.

\subsection{Why the projection theorem matters}

\Cref{thm:projection} gives a concrete target for future MTT work.  If the
selected upper action is found, one need not compare final amplitudes by
analogy.  One can verify the commuting projection at the level of the
quadratic operator, every interaction tensor, grading, and identity
operator.  Equality of the retained graph coefficients then follows
mechanically.

The theorem also identifies failure modes. If \(KP\ne P^*K\),
the lower propagator is not simply \(PGP^*\). If mixed vertices
survive, discarded modes generate effective interactions.  If BRST/BV
operators fail to intertwine, gauge identities need not descend.  These
are informative obstructions, not merely missing prose.

\subsection{Best next theorem}

The decisive upstream theorem is:
\[
\begin{gathered}
\text{selected q79 upper action and complex}\\
\Downarrow\\
K,\ \{V_n\},\ Q_{\BRST},\ \text{state/renormalization data}\\
\Downarrow\quad\text{by certified intertwiners}\\
\text{accepted lower free and interacting structures}.
\end{gathered}
\]
This would turn the graph datum from an input into a derived object.  Until
then, the present paper supplies the exact transfer mathematics and the
boundary that any claimed derivation must cross.

\section{Conclusion}

The recurrence of Feynman graphs has a clean explanation.  A quadratic
contraction and multilinear interactions generate graph-indexed
perturbative coefficients through graded Wick expansion.  This grammar is
shared by classical and quantum models.  It is not, on its own, a
quantization theorem.

We proved that compatible finite projection preserves every all-retained
graph coefficient and, under absence of mixed vertices, the normalized
retained perturbative expansion.  We also proved that graph combinatorics
cannot select numerical weights, statistics, causal prescriptions, gauge
data, or renormalization.  These statements give modal diagrammatics a
precise role: it is an exact language for transferring a \emph{selected}
perturbative datum, not a substitute for selecting that datum.

For MTT, the free q79 even-CAR net is already available through a companion
construction using standard CAR/AQFT machinery.  The next frontier is the
same-source upper action and its interacting gauge--BRST completion.  A
future success there can plug directly into the theorem proved here.

% BEGIN MTT MANAGED COMPUTATIONAL EVIDENCE
\section*{Computational Evidence and Reproducibility}
The theorems proved here are finite algebraic statements and do not depend
on numerical fitting.  The curated repository
\url{https://github.com/PeterNero/mtt-results-repro} is cited for current
MTT status context.  In particular,
\path{A04/final_12_of_12_audit} records embedded
renormalized-Standard-Model equivalence at the adopted profile tier, while
\path{A05/strict_upgrade_ledger} records stronger no-knob obligations
that remain open.  Neither row proves
\cref{thm:wick-graph,thm:projection,thm:nonselection}, and the open row is
not evidence of closure.

The selected free-CAR theorem and the classical-to-quantum interface are
documented in the companion records
\cite{MTTQFT2026,MTTBasinQFT2026}.  Exact release artifacts, source hashes,
and the revision audit accompany this paper in the MTT papers repository.
% END MTT MANAGED COMPUTATIONAL EVIDENCE

\bibliographystyle{plain}
\bibliography{main}

\end{document}
